\SourceLeaf{263}%
too small to detect. For example, in the largest triangle we measured in previous years, the one between Hohenhagen, Brocken and Inselsberg, the sum of the angles had an excess \(=14''85348\). The calculation gave the following corrections to the angles.
\begin{center}
\begin{tabular}{ll}
Hohenhagen \(.\ .\ .\ .\ .\)&\(-4''95113\)\\
Brocken \(.\ .\ .\ .\ .\ .\)&\(-4{,}95104\)\\
Inselsberg\(.\ .\ .\ .\ .\ .\)&\(-4{,}95131\)
\end{tabular}
\end{center}
\Article{29}
\indent To conclude, we shall also compare the area of a triangle on a curved surface with that of a rectilinear triangle whose sides are \(a,b,c\). We denote the latter area by \(\sigma^*\) \(=\tfrac{1}{2}bc\sin A^*=\tfrac{1}{2}ac\sin B^*=\tfrac{1}{2}ab\sin C^*\)

Up to quantities of the fourth order, we have
\[
\sin A^*=\sin A-\tfrac{1}{12}\sigma\cos A\mathord{.}(2\alpha+\beta+\gamma)
\]
or, to the same accuracy,
\[
\sin A=\sin A^*\mathord{.}(1+\tfrac{1}{24}bc\cos A\mathord{.}(2\alpha+\beta+\gamma))
\]
Substituting this value in formula [9] gives, up to quantities of the sixth order,
\[
\begin{aligned}
\sigma=\tfrac{1}{2}bc\sin A^*\mathord{.}\{1&+\tfrac{1}{120}\alpha(3bb+3cc-2bc\cos A)\\
&+\tfrac{1}{120}\beta(3bb+4cc-4bc\cos A)\\
&+\tfrac{1}{120}\gamma(4bb+3cc-4bc\cos A)\}
\end{aligned}
\]
or, to the same accuracy,
\[
\begin{aligned}
\sigma=\sigma^*\{1&+\tfrac{1}{120}\alpha(aa+2bb+2cc)+\tfrac{1}{120}\beta(2aa+bb+2cc)\\
&+\tfrac{1}{120}\gamma(2aa+2bb+cc)\}
\end{aligned}
\]
For a spherical surface this formula takes the form
\[
\sigma=\sigma^*(1+\tfrac{1}{24}\alpha(aa+bb+cc))
\]
It is easy to verify that the following formula can be used in its place with the same precision
\[
\sigma=\sigma^*\BareRoot{\frac{\sin A\mathord{.}\sin B\mathord{.}\sin C}{\sin A^*\mathord{.}\sin B^*\mathord{.}\sin C^*}}
\]
If the same formula is applied to triangles on a curved surface that is not spherical, the error will generally be of the fifth order. It will be too small to detect in any triangle of the kind that can be measured on the Earth’s surface.
\par\medskip\centerline{\rule{.75in}{.3pt}}
